%=======================================================================
%  Paper V — Section 4 (The four weak-signal bodies)   FINAL
%
%  Requires the macros listed at the head of Section 2.
%  Additional macro used here:
%      \newcommand{\densregion}{\varrho_{R}}   % already defined as \densR
%
%  This version supersedes all earlier drafts of Section 4. Principal
%  changes against ver10:
%      contrast intervals    -> now include the induced change in
%                               rho_rest, as in Eq. (itokawaref)
%      Eros "1.000 +/- 0.017"-> 1.018 +/- 0.026 (central value corrected)
%      Bennu I = 0.000 %     -> I < 10^-3 %, with the floor stated
%      negative I values     -> identified as a grid artefact and bounded
%      amplification factor  -> tabulated; ranks the bodies identically
%                               to I (Spearman rho = -1 over five bodies)
%      S statistic           -> shown to FAIL on the Eros misplaced plane,
%                               and demoted to a necessary-not-sufficient
%                               condition within a three-part criterion
%
%  All values confirmed against paperV_final.json, paper4_recheck.json,
%  v9_section2_tests.json and arro_235/250/500.json, arrokoth_clean.json.
%  Corrections in this pass: S at the misplaced Eros plane 19 -> 2.8
%  (the region holds 88 per cent of the volume, not 5-20); Arrokoth run at
%  235 kg/m3 added, reversing the sign; Arrokoth table column count fixed.
%=======================================================================

\section{The four weak-signal bodies}
\label{sec:nulls}

Itokawa is one body out of five. This section reports the other four, and
its purpose is not to add results but to bound a failure rate: if the
inversion manufactured structure wherever it was pointed, the method would
be worthless whatever it returned on Itokawa. It does not. All four return
improvements at or below the mesh-resolution sensitivity envelope of
Paper~IV, and all four return contrasts consistent with unity once the
curvature sensitivity interval is propagated through the mass constraint.

Three further things are established here, and they matter more than the
null results themselves. The amplification factor $\ampfac$ of
Equation~\eqref{eq:amplification}, computed from densities alone and
without reference to the objective function, ranks all five bodies in
exactly the order of their dispersion improvements
(Section~\ref{sec:nulls:amplification}). A deliberately misplaced plane on
Eros produces a displacement statistic $\Sstat$ three times larger than
Itokawa's, which shows that $\Sstat$ cannot be used on its own
(Section~\ref{sec:nulls:misplaced}). And on Arrokoth the recovered
contrast reverses sign along two independent axes---the assumed bulk
density and the position of the dividing plane---which is the concrete
form of the statement that nothing is invariant when the signal is weak
(Section~\ref{sec:nulls:arrokoth}).

\subsection{The four results}
\label{sec:nulls:table}

Table~\ref{tab:nulls} collects them. Two columns require comment before
the bodies are taken individually.

The contrast intervals are computed as in
Equation~\eqref{eq:itokawaref}: the curvature sensitivity interval on
$\densR$ and the change it induces in $\densrest$ through the mass
constraint move the ratio in the same direction, so their fractional
contributions add rather than combining in quadrature,
%
\begin{equation}
\frac{\delta C}{C}
 = \frac{\delta\densR}{\densR}
 + \frac{f}{1-f}\,\frac{\delta\densR}{\densrest} .
\label{eq:contrastinterval}
\end{equation}
%
On Bennu, where the region is half the body, the second term is as large
as the first and the interval is twice what the first term alone would
give. Reporting only the first term, as earlier drafts of this section
did, understates every interval in the table.

The improvement for Bennu is quoted as a bound rather than as a value.
Because no sub-grid refinement is applied
(Section~\ref{sec:methods:minimisation}), the scan need not contain the
density at which $\potvarn$ is stationary, and the reported minimum can
therefore lie marginally \emph{above} the uniform solution, giving a
negative $\improve$. This is an arithmetic impossibility for the
underlying problem---the uniform body is the special case
$\densR = \densrest$ of the two-region model, so the minimum over
$\densR$ can never exceed it---and its appearance is a measure of the
grid rather than of the physics. It is observed on Arrokoth
(Section~\ref{sec:nulls:arrokoth}), where it reaches
$\improve = -0.006\%$, and that figure is the empirical floor of the
calculation. We therefore report improvements below $10^{-3}\%$ as
bounds and do not distinguish among them.

\begin{table}[htbp]
\centering
\small
\setlength{\tabcolsep}{4pt}
\caption{The four weak-signal bodies. Planes are those of
Table~\ref{tab:planes}; ``imposed'' marks the one plane not returned by
the detector. $f = \VR/\Vtot$. Contrast intervals follow
Equation~\eqref{eq:contrastinterval} and are therefore wider than the
curvature interval on $\densR$ alone. $\ampfac$ is the amplification
factor of Equation~\eqref{eq:amplification}, computed from the densities
without reference to the objective; its sign follows that of
$\densbulk - \densrest$ and is negative where the region is the less
dense of the two. $\Sstat$ is the displacement statistic of
Equation~\eqref{eq:sstat}. Itokawa is repeated from
Section~\ref{sec:itokawa} for comparison and is not part of this section's
argument.}
\label{tab:nulls}
\setlength{\tabcolsep}{10pt}
\begin{tabular}{llrrrrrr}
\toprule
Body & Plane & $f$ & $\densR$ & $\densrest$ & $C = \densR/\densrest$
 & $\improve$ & $\dCOM$ \\
 & (km) & (\%) & \multicolumn{2}{c}{(kg\,m$^{-3}$)} & & (\%) & (m) \\
\midrule
Bennu    & $0$ (imposed)   & $50.2$ & $1190$ & $1192$   & $0.999\pm0.071$ & $<10^{-3}$ & $-0.07$ \\
Eros     & $+4.4698$       & $34.6$ & $2700$ & $2652$   & $1.018\pm0.026$ & $0.418$    & $+56.4$ \\
67P      & $+0.9517$       & $26.6$ & $565$  & $527$    & $1.073\pm0.075$ & $0.574$ & $+29.1$ \\
Arrokoth & $-4.7863$       & $66.3$ & $390$  & $419.7$  & $0.929\pm0.038$ & $1.272$ & $-141$ \\
\midrule
\itshape Itokawa & \itshape $+0.1608$ & \itshape $16.4$ & \itshape $2800$
 & \itshape $1858$ & \itshape $1.51\pm0.09$ & \itshape $16.9$ & \itshape $+16.9$ \\
\bottomrule
\end{tabular}

\vspace{4pt}
\begin{minipage}{\textwidth}
\footnotesize
Arrokoth is quoted at an assumed bulk density of
$\SI{400}{\kilo\gram\per\cubic\metre}$; the dependence on that assumption
is the subject of Section~\ref{sec:nulls:arrokoth}. Its contrast interval
uses the curvature interval of the run with the buried facets retained,
$\pm\SI{5.6}{\kilo\gram\per\cubic\metre}$, the rerun that excludes them
having been made without it; the two runs differ by one scan step in
$\densR$. The 67P curvature interval is $\pm\SI{28}{\kilo\gram\per\cubic\metre}$.
\end{minipage}
\end{table}

\subsection{The amplification factor ranks the bodies}
\label{sec:nulls:amplification}

A result emerges from Table~\ref{tab:nulls} that we did not anticipate
and that costs nothing to obtain. The amplification factor
$\ampfac = \densrest/(\densbulk-\densrest)$ is computed from two
densities. It involves no surface average, no facet areas, no rotation
rate, and no evaluation of $\potvarn$; given a shape model, a mass and a
dividing plane, it is available before the objective function is ever
called. Yet it orders the five bodies exactly as the dispersion
improvement does.

\begin{table}[htbp]
\centering
\small
\caption{The amplification factor against the dispersion improvement. The
two orderings are identical, with $\ampfac$ increasing as $\improve$
decreases; over the five bodies, each at the plane its
detector returned, the Spearman rank correlation is $-1.000$ exactly
($p = 0.017$, two-tailed); the ordering does not survive the inclusion of a
misplaced plane, as the text below records. $\ampfac$ diverges as the
signal vanishes, because its denominator $\densbulk - \densrest$ is
proportional to the excess mass and the excess mass goes to zero.}
\label{tab:amplification}
\setlength{\tabcolsep}{10pt}
\begin{tabular}{lrrr}
\toprule
Body & $\densbulk-\densrest$ (kg\,m$^{-3}$) & $\lvert\ampfac\rvert$
 & $\improve$ (\%) \\
\midrule
Itokawa  & $+154.8$ & $12.0$   & $16.9$ \\
Arrokoth & $-19.7$  & $21.3$   & $1.272$ \\
67P      & $+10.1$  & $52.2$   & $0.574$ \\
Eros     & $+15.8$  & $168$    & $0.418$ \\
Bennu    & $-1.0$   & $1187$   & $<10^{-3}$ \\
\bottomrule
\end{tabular}
\end{table}

The reason is not mysterious once
Equation~\eqref{eq:excessidentity} is in view. Both quantities are
monotone functions of the excess mass: $\improve$ measures how much
dispersion the excess mass removes, and $\ampfac^{-1}$ measures how large
the excess mass is relative to the body. What is useful is that the
second can be computed without the first, so that $\ampfac$ serves as a
diagnostic on any body for which a mass and a plane are available---and
that a large $\ampfac$ is a warning, not a precision. A body returning
$\ampfac = 1187$, as Bennu does, is one on which the two-region model has
found essentially no excess mass, and on which any fractional stability
quoted for $\densrest$ corresponds to a fractional instability in
$\Dmass$ three orders of magnitude larger.

We report this as a correlation over five bodies and not as a law. Five
points, four of them nulls, cannot establish a functional form, and the
ordering could be coincidental at the $1.7\%$ level. It is offered as a
cheap first check, and Section~\ref{sec:guidelines} recommends it as one.

One limit on that check is already visible in this paper and should be
stated here rather than left to be discovered. The five bodies of
Table~\ref{tab:amplification} are evaluated at planes chosen by the rule of
Section~\ref{sec:methods:neck}. At the misplaced Eros plane of
Section~\ref{sec:nulls:misplaced} the same formula gives
$\densrest = \SI{2996.9}{\kilo\gram\per\cubic\metre}$ against a bulk
density of $2667.8$, so $\ampfac = -9.11$: in absolute value smaller than
Itokawa's $12.0$, which on the ordering above would rank a spurious result
as the strongest signal in the study. The excess mass it implies,
$-12.3\%$ of the total, is larger than Itokawa's $+7.7\%$ for the same
reason. $\ampfac$ therefore orders bodies whose planes are properly placed,
and says nothing about whether a plane is properly placed; used without the
provenance checks of Section~\ref{sec:failures:plane} it will rank a
misplaced plane above a real signal. The displacement statistic behaves
better on this case, returning $2.86$ against Itokawa's $5.8$
(Section~\ref{sec:failures:plane}), but it does not reject it either.

\subsection{Eros at the detected saddle}
\label{sec:nulls:eros}

The detector returns the Himeros saddle at $x = \SI{+4.4698}{\kilo\metre}$
with a prominence of $0.216$, the second highest in the study, placing
$34.6\%$ of the volume in the region. The inversion returns
$\densR = 2700$ against $\densrest = \SI{2652}{\kilo\gram\per\cubic\metre}$,
a contrast of $1.018$, and the dispersion falls from $0.03413$ to
$0.03399$, an improvement of $0.418\%$.

The curvature sensitivity interval is
$\pm\SI{44}{\kilo\gram\per\cubic\metre}$ on $\densR$, which through
Equation~\eqref{eq:contrastinterval} gives $C = 1.018 \pm 0.026$: the
contrast is $0.7$ standard intervals from unity and consistent with a
homogeneous body. We state the central value as $1.018$ rather than as
unity, because the two are not the same statement and the first is what
the inversion returned.

Three figures place the result. The improvement of $0.418\%$ lies below
the whole of the $0.9$--$5.9\%$ envelope, so no comparison with the
envelope's edges is needed. The amplification factor is $\ampfac = 168$,
fourteen times the Itokawa value. And the displacement statistic is
$\Sstat = 0.68$: the background density sits less than one curvature
interval from the bulk density, which is what a uniform body does.

One caution attaches to the improvement. The absolute change in
$\potvarn$ is $1.4\times10^{-4}$, and the field-point offset of
Section~\ref{sec:methods:model} perturbs the surface potential on Eros by
of order $0.04\%$ of $\lvert\Ubar\rvert$. Whether that perturbation
propagates into $\potvarn$ at the $10^{-4}$ level is exactly what the
offset-sensitivity test of that section measures, and the Eros null
should be read together with its outcome rather than before it. That
test returns $0.426\%$ with the offset scaled by $10$ and $0.418\%$ with it
scaled by $0.1$, against $0.419\%$ at the nominal offset: the improvement
moves by $0.007$ percentage points, and the recovered contrast not at all,
so the Eros null does not depend on where the field points sit.

\subsection{Eros with the plane misplaced, and the limits of \texorpdfstring{$\Sstat$}{S}}
\label{sec:nulls:misplaced}

Section~\ref{sec:methods:neck} rejected the rule that selects the
smallest interior half-width, on the grounds that on a body tapering at
both ends it lands on a flank. Eros is such a body, and the rejected rule
selects $x = \SI{-10.7619}{\kilo\metre}$. Running the inversion there is
the one worked failure case of this paper, and it is more instructive
than we expected.

At that plane the inversion returns a contrast of $0.876$---a density
deficit of over twelve per cent, of the opposite sign to the Itokawa
result and larger in magnitude than three of the four nulls---with an
improvement of $\improve = 5.526\%$ and a centre-of-mass offset of
$\SI{-214.8}{\metre}$, nearly four times the offset at the true saddle.
Nothing in the output marks it as spurious. It would be reported as a
detection by any criterion based on the size of the recovered contrast.

The improvement is the first thing that catches it, and it does so by a
narrow margin that we state rather than round away. At $5.526\%$ it falls
inside the $0.9$--$5.9\%$ envelope, but only $6\%$ below the upper edge,
and that upper edge is attained on a different body. The comparison that
matters is against the value the envelope takes on Eros itself, which is
$6.03\%$ over the ladder from $2{,}000$ to $42{,}000$ facets with quadric
decimation; the misplaced plane is $0.92$ times that value, and the test
rejects it with a margin of eight per cent. The margin depends on the
decimator. With vertex clustering, which Paper~IV used, the Eros envelope is
$5.40\%$, the misplaced plane is $1.02$ times it, and the test would
\emph{not} reject the case. We use the quadric value because it is the one
consistent with the meshes of this paper, but a reader should know that the
envelope test catches this artefact by a margin smaller than the difference
between two defensible ways of computing the envelope.

The second thing that catches it is provenance: the plane was selected by
a rule the paper had already rejected on geometric grounds, before any
density was computed. That is a strong argument, but it is not available
to a reader who is handed only the output.

The displacement statistic ranks it correctly but does not reject it. The
curvature sensitivity interval at the misplaced plane is
$\pm\SI{15}{\kilo\gram\per\cubic\metre}$, three times \emph{narrower}
than at the true saddle: the objective is more sharply curved about the
spurious minimum than about the real one. Propagating it through the mass
constraint,
%
\begin{equation}
\Sstat = \frac{\lvert\densbulk-\densrest\rvert}{\sigma(\densrest)}
       = \frac{\densbulk\,(1-C)f/[1-(1-C)f]}
              {[f/(1-f)]\,\delta\densR}
       = \frac{329}{116} = 2.8 ,
\label{eq:sfail}
\end{equation}
%
evaluated with the region holding $f = 0.884$ of the volume, $C = 0.876$
and $\densbulk = \SI{2668}{\kilo\gram\per\cubic\metre}$. The misplaced
plane leaves the region on the large side of the cut, so the small
complement makes $\sigma(\densrest)$ some eight times $\delta\densR$, and
that factor, not the sharpness of the minimum, sets the value. $\Sstat = 2.8$
is half the $5.8$ returned by the Itokawa reference solution, so the
statistic does rank the artefact below the genuine result. It does not
reject it: a displacement of nearly three curvature intervals is not one a
reader handed only this number would dismiss.

We draw the conclusion rather than leave it implicit. $\Sstat$ measures
displacement in units of curvature, and a sharply curved objective makes
a small displacement look significant; sharpness of the minimum is not
evidence that the minimum is in the right place. $\Sstat$ is therefore a
necessary condition and not a sufficient one, and
Sections~\ref{sec:methods:minimisation} and~\ref{sec:itokawa:ambiguity},
which use it, should be read with that qualification. We report a result
as a detection only where all three of the following hold:
%
\begin{enumerate}
\item the improvement exceeds the value the mesh-resolution envelope
takes on that body;
\item the displacement statistic $\Sstat$ is large, in the sense of
Equation~\eqref{eq:sstat}; and
\item the dividing plane carries geometric warrant---a saddle returned by
the detector of Section~\ref{sec:methods:neck} at a prominence above
threshold, rather than a plane chosen by the investigator.
\end{enumerate}
%
The misplaced Eros plane passes the second and fails the first and third.
Itokawa passes all three. No other combination arises in this study, and
we do not claim the criterion is complete; it is the weakest set of
conditions that rejects every case we have shown to be spurious.

\subsection{Bennu: a structure the model cannot represent}
\label{sec:nulls:bennu}

Bennu is the one body whose plane we imposed. The detector refuses it at
a prominence of $0.0004$, more than two orders of magnitude below the
lowest accepted value, and correctly: Bennu is an oblate spheroid with no
constriction. We placed a plane through the centre of figure, which
divides the volume $50.2$ to $49.8$, precisely to see what the objective
returns when the model has no geometric warrant at all.

It returns nothing. The recovered densities are $\densR = 1190$ and
$\densrest = \SI{1192}{\kilo\gram\per\cubic\metre}$, a contrast of
$0.999$; the curvature interval of
$\pm\SI{42}{\kilo\gram\per\cubic\metre}$ gives $C = 0.999 \pm 0.071$
through Equation~\eqref{eq:contrastinterval}, the widest interval in the
study because the region is half the body and the second term of that
equation is as large as the first. The improvement is below the
$10^{-3}\%$ floor established in Section~\ref{sec:nulls:table}. The
centre of mass moves by $\SI{7}{\centi\metre}$ on a body of mean radius
$\SI{245}{\metre}$. The amplification factor is $1187$ and the
displacement statistic is $\Sstat = 0.02$.

This is the cleanest null in the study, and it is worth being precise
about what it does and does not show. It does not show that Bennu is
homogeneous. Radio science and the OSIRIS-REx gravity field indicate the
opposite: an underdense equatorial bulge and an underdense core
\citep{scheeres2020}. A plane-bounded two-region model is blind to both,
because a radial or axisymmetric density variation contributes almost
nothing to the difference between the two half-spaces on either side of a
plane through the centre---the deficit is distributed symmetrically about
it. The return of unity is the correct response of a model to a structure
it cannot express, not a measurement of the structure.

Read the other way, this is the strongest available bound on the
false-positive rate of the method on a real body: given a real shape
model, a real mass, a real spin state, and a plane with no geometric
justification whatever, the inversion declined to invent a contrast. The
eight synthetic bodies of Section~\ref{sec:constrains} extend that bound;
Bennu anchors it.

\subsection{67P: a signal below every systematic}
\label{sec:nulls:67p}

The detector returns the neck at $x = \SI{+0.9517}{\kilo\metre}$ with a
prominence of $0.091$, the lowest accepted in the study, placing $26.6\%$
of the volume in the small lobe---the figure validated against
\citet{jorda2016} in Section~\ref{sec:methods:external}. The inversion
returns $\densR = 565$ against $\densrest = \SI{527}{\kilo\gram\per
\cubic\metre}$, a contrast of $1.073$, with $\improve = 0.574\%$.

The number is inside the envelope, and that alone would settle it. But on
67P a stronger statement is available, because the systematic
uncertainties on this body are all larger than the signal. Repeating the
inversion at the same plane on the second of the two acceptable shape
models returns $\densrest = \SI{521}{\kilo\gram\per\cubic\metre}$ and
$\improve = 0.725\%$: the same $\densR$, but an improvement $26\%$ larger
in relative terms. The decimation from $1.3$ million facets to $50{,}000$
moves the recovered contrast by $-3.9\%$
(Section~\ref{sec:methods:shapes}), a figure some seven times the
improvement itself. And the choice of shape model changes $\Vtot$ by
$1.1\%$ between SPG and SPC products. A signal smaller than the
model-to-model scatter, smaller than the decimation error, and smaller
than the volume difference between reductions of the same data is not a
signal.
% Confirmed: Table 4.1 reports the MSPCD model, a closed surface (chi = 2)
% after decimation. The SPC 2017 model, not closed as supplied, enters only
% the cross-model comparison and no tabulated result.

The neck model is the one place in this section where an alternative
region shape does better. A slab spanning the constriction returns a
contrast of $0.903$---a \emph{deficit}, the neck less dense than the
lobes---with $\improve = 1.132\%$, twice the head-model value. Both lie
inside the envelope and neither is a detection, but the ordering is worth
recording because it is the reverse of Itokawa's, where the head model
beat the neck model by a factor of seven
(Section~\ref{sec:itokawa:ambiguity}), and it is of the opposite sign,
Itokawa's neck model having returned a contrast of $1.300$. Which region
shape the objective prefers is a property of the individual body, not a
general feature of contact binaries, and a study that tested only one
shape would not learn this.

A lower-density neck is what a compressed contact between two lobes would
not produce and what a region of enhanced porosity would; the point is
raised only to note that the sign is physically interpretable and that
the magnitude does not support interpreting it.

\subsection{Arrokoth: two independent sign reversals}
\label{sec:nulls:arrokoth}

Arrokoth returns the largest improvement of the four, $1.272\%$, and is
the only body in the study without a measured mass. The two facts are
connected, and the connection is the subject of this subsection.

\paragraph{Dependence on the assumed bulk density.}
With no mass determination \citep{spencer2020}, $\Mtot$ must be assumed,
and the assumption propagates directly into the recovered contrast.
Table~\ref{tab:arrokothdensity} gives the four values run, and
Figure~\ref{fig:arrokoth} plots them.

\begin{figure}[htbp]
\centering
\includegraphics[width=0.72\textwidth]{fig05_arrokoth_mass.pdf}
\caption{Recovered contrast on Arrokoth against the assumed bulk density,
with and without the buried facets of
Section~\ref{sec:methods:shapes}. Excluding them, the contrast crosses unity
between $235$ and $\SI{250}{\kilo\gram\per\cubic\metre}$: at the central
published estimate of \citet{keane2022} the region is the denser of the two,
and at every higher assumption it is the lighter. Retaining them shifts the
whole curve upward and moves the crossing to about
$\SI{345}{\kilo\gram\per\cubic\metre}$. The shaded region is excluded by the
lower bound near $\SI{290}{\kilo\gram\per\cubic\metre}$ that
\citet{spencer2020} derive from the requirement that mutual gravity exceed
the centrifugal acceleration; the crossing lies just below it.}
\label{fig:arrokoth}
\end{figure}

\begin{table}[htbp]
\centering
\small
\caption{Arrokoth at four assumed bulk densities, at the detected plane
$x = \SI{-4.7863}{\kilo\metre}$ with the buried facets of
Section~\ref{sec:methods:shapes} excluded. The recovered contrast rises
toward unity as the assumed density falls, reaches it at $250$, and crosses
it: at $235$ the region is the denser of the two. The improvement passes
through the negative value that marks the grid floor of
Section~\ref{sec:nulls:table} at the crossing.}
\label{tab:arrokothdensity}
\setlength{\tabcolsep}{10pt}
\begin{tabular}{rrrrrr}
\toprule
Assumed $\densbulk$ & $\densR$ & $\densrest$ & $C$ & $\Dmass/\Mtot$ & $\improve$ \\
(kg\,m$^{-3}$) & \multicolumn{2}{c}{(kg\,m$^{-3}$)} & & (\%) & (\%) \\
\midrule
$235$ & $240.0$ & $225.2$ & $1.0659$ & $+4.19$ & $+1.224$ \\
$250$ & $250.0$ & $250.0$ & $0.9999$ & $-0.01$ & $-0.006$ \\
$400$ & $390.0$ & $419.7$ & $0.9293$ & $-4.92$ & $+1.272$ \\
$500$ & $485.0$ & $529.5$ & $0.9159$ & $-5.91$ & $+2.378$ \\
\bottomrule
\end{tabular}
\end{table}

The published bulk density is not well constrained, and the constraints do
not agree. \citet{keane2022} infer $\SI{235}{\kilo\gram\per\cubic\metre}$
with a one-sigma range of $155$ to $\SI{600}{\kilo\gram\per\cubic\metre}$;
\citet{spencer2020} require more than $290$ if the neck is not in tension;
\citet{mckinnon2020} find $250$ to $500$ for a comet-like strength. We ran the
central value of \citet{keane2022} as well as three values spanning the
range of \citet{mckinnon2020}.

At $\SI{235}{\kilo\gram\per\cubic\metre}$ the recovered contrast is
$1.066$, the region the denser; at $250$ it is $0.9999$; at $400$ and $500$ it
is $0.929$ and $0.916$, the region the less dense. The sense of the answer
reverses between $235$ and $250$, inside the published one-sigma range. A
linear extrapolation from the three higher values would have put the
contrast at $235$ near $1.006$; the measured value is ten times further from
unity, so the crossing could not have been located from runs on one side of
it. The inversion returns a surplus of seven per cent at the density
Arrokoth most probably has, and a deficit of six to eight per cent at
densities near the upper end of the published range. What the method
reports about this body is, to a good approximation, a restatement of what
was assumed about it.

\paragraph{Dependence on the buried facets.}
The Arrokoth model is supplied as two unjoined shells, and $455$ facets,
carrying $0.943\%$ of the surface area, lie inside the opposite shell
(Section~\ref{sec:methods:shapes}). They are not part of the body's
surface and the relaxation argument of \citet{richardson2014} does not
apply to them, so all figures above exclude them. With them
excluded the contrast crosses unity between $235$ and
$\SI{250}{\kilo\gram\per\cubic\metre}$, and linear interpolation between
the two runs puts the crossing within a kilogram per cubic metre of $250$.
Retaining them moves the crossing to approximately
$\SI{345}{\kilo\gram\per\cubic\metre}$---a shift of some $38\%$,
produced by facets carrying under one per cent of the area. At $235$ the two
treatments agree, both returning $1.066$, so the buried facets move the
crossing without changing the answer far from it.
A sub-percent geometric decision therefore controls the sign of the
result over a wide band of assumed densities. This is why
Section~\ref{sec:methods:model} recomputes $\Ubar$ and $\sum_i A_i$ over
the retained set rather than merely omitting terms from the numerator.

\paragraph{Dependence on the plane.}
The third dependence is the one that makes this body decisive for the
argument of Section~\ref{sec:constrains}. Sweeping the plane so that the
region fraction runs from $40.0\%$ to $75.6\%$, at a fixed assumed bulk
density of $\SI{400}{\kilo\gram\per\cubic\metre}$, the recovered density
difference $\Ddens$ runs from $+25.0$ to
$\SI{-61.5}{\kilo\gram\per\cubic\metre}$. It passes through zero at
$f = 61.2\%$, where the inversion returns
$\densR = \densrest = \SI{400}{\kilo\gram\per\cubic\metre}$ exactly, a
contrast of $1.0000$, and an improvement at the negative grid floor.

The consequence for the excess mass is the point. Since
$\Dmass/\Mtot = \Ddens f/\densbulk$, the sweep carries $\Dmass$ from
$+2.5\%$ of the total to $-11.6\%$: a swing of fourteen percentage points
that changes sign in the middle. On Itokawa the corresponding quantity
varied by $\pm5.12\%$ about a non-zero mean and never approached zero
(Table~\ref{tab:itokawasweep}). Here there is no invariant to speak of,
because the quantity that was invariant on Itokawa is identically zero
somewhere inside the sweep, and a quantity that passes through zero has
no meaningful fractional stability.

The detected plane, at $f = 66.3\%$, lies $5.1$ percentage points from
that null. The contrast of $0.929$ reported in Table~\ref{tab:nulls} is
therefore principally a measure of how far the detector's plane happens
to fall from the point at which the model degenerates, and not of any
property of the body. This is the clearest demonstration in the paper of
the claim made in the abstract: where the signal is weak, nothing is
invariant, and the apparent stability documented on Itokawa is a
property of the strong-signal regime rather than of the method.

\paragraph{Resolution and the neck model.}
Two further checks complete the body. Repeating the inversion on meshes
of $15{,}000$, $25{,}000$ and $40{,}960$ facets returns
$\densR = \SI{395}{\kilo\gram\per\cubic\metre}$ on all three, with the
complementary density between $409.8$ and
$\SI{410.2}{\kilo\gram\per\cubic\metre}$: a variation of
$\SI{0.4}{\kilo\gram\per\cubic\metre}$, or $0.1\%$, and, unlike the
Itokawa case of Section~\ref{sec:itokawa:mesh}, with the detected plane
displaced by no more than one scan step across the three meshes.
This series retains the buried facets, and that alone accounts for the
difference from Table~\ref{tab:arrokothdensity}. The plane is the same, and
so is the volume fraction: $66.330\%$, $66.328\%$ and $66.326\%$ on the three
meshes against $66.326\%$ in the table, the three-thousandth of a per cent
being the volume the decimation loses. Both answers satisfy the mass
constraint at that fraction---$0.66326\times395 + 0.33674\times409.8 = 400.0$
and $0.66326\times390 + 0.33674\times419.7 = 400.0$---because $\densrest$ is
solved from it rather than fitted, so the constraint cannot distinguish them.
What distinguishes them is the surface over which the dispersion is averaged:
with the $455$ buried facets included the minimum falls one scan step higher,
at $395$ rather than $\SI{390}{\kilo\gram\per\cubic\metre}$, and the contrast
reads $0.963$ rather than $0.929$. They are two objective functions evaluated
on the same body at the same plane, not two answers to the same question. We
keep the buried facets in this series because it compares meshes with one
another, and identifying the buried set separately on each decimated mesh
would introduce a second difference between the runs.

The neck model does not survive at all. Its minimum falls on the lower
limit of the scan interval, $\SI{25}{\kilo\gram\per\cubic\metre}$, with a
slab enclosing $1.0\%$ of the volume: the small-region limit in which the
mass constraint forces an arbitrarily large density contrast for an
arbitrarily small mass, and the objective is driven by the scan boundary
rather than by the data. The boundary-minimum guard of
Section~\ref{sec:methods:minimisation} flagged it and the result is
discarded. It is one of two occasions on which that guard fired: the
other, at $\xsplit = \SI{0.2200}{\kilo\metre}$ on Itokawa, was resolved
by widening the interval and is described in
Section~\ref{sec:failures}. Here no widening helps, because the limit is
physical---a slab of one per cent of the volume cannot carry a meaningful
mass anomaly---and the model is simply inapplicable.

\subsection{What the four cases establish}
\label{sec:nulls:summary}

The nulls are not symmetrical with the detection, and it is worth being
explicit about the asymmetry. Four bodies returning nothing bound the
rate at which this method invents structure; they do not bound the rate
at which it misses structure, because on none of them is the true
interior known. Bennu is the sharpest illustration: the inversion
returned unity on a body that independent gravity data show to be
heterogeneous, and that is a false negative by construction, deliberately
provoked.

What the four do establish is the following. The method does not
manufacture a contrast when handed a plane with no geometric warrant
(Bennu). It does not manufacture one on a body whose neck is real but
whose interior contrast, if any, is below the systematics of the
available shape models (67P). It returns an improvement below the
resolution envelope on a body with the second-highest neck prominence in
the study (Eros). And on a body without a measured mass it returns a
result whose sign is controlled by three separate modelling choices, none
of which is an observation (Arrokoth).

Against this, one case shows that the method \emph{can} be made to
produce a plausible-looking contrast: the misplaced plane on Eros, which
yields a twelve per cent density deficit, an improvement just inside the
envelope, and a displacement statistic three times the value returned by
the genuine Itokawa result. That case is rejected here on the provenance
of its plane and on the envelope comparison, not on any property of its
output, and it is the reason Section~\ref{sec:failures} catalogues
failure modes rather than merely listing null results.